\chapter{系统测试}

\section{测试环境与工具}

为验证系统设计与实现的正确性，本文围绕认证授权、订单处理、库存管理、生产流转、采购协同、质量追溯与实时通知等功能开展测试。测试环境覆盖后端运行环境、数据库环境和前端访问环境，具体配置如表~\ref{tab:test-env}所示。

\begin{table}[htbp]
\centering
\footnotesize
\caption{系统测试环境}
\label{tab:test-env}
\begin{tabular}{|c|c|c|}
\hline
类别 & 配置项 & 说明 \\\hline
操作系统 & Windows & 系统部署与测试运行环境 \\\hline
JDK & Java 17 & 后端运行环境 \\\hline
构建工具 & Maven、Node.js & 前后端依赖安装与项目构建 \\\hline
数据库 & MySQL & 业务数据存储与测试数据准备 \\\hline
前端运行环境 & Vue 3、Vite & 页面访问与交互验证 \\\hline
浏览器 & Chrome 或 Edge & 功能页面与实时通知验证 \\\hline
\end{tabular}
\end{table}

测试过程中采用模块化验证思路：先对关键服务进行单元级检查，再对跨角色流程进行功能测试，以确保系统既满足局部业务规则，又能够在整体流程上保持贯通。

\section{单元测试}

单元测试主要面向后端关键服务层逻辑，重点验证认证模块、订单流程、生产流程、采购流程、质检流程以及周计划生成服务等是否符合预期。测试内容主要包括用户登录后令牌生成与解析是否正确、订单状态转换是否符合约束、库存更新是否与库存流水一致、生产完成后是否正确进入待质检状态、采购订单状态变化是否准确，以及周计划生成结果是否能够正确识别缺口并给出建议。

\begin{table}[htbp]
\centering
\footnotesize
\caption{关键单元测试内容}
\label{tab:unit-test}
\begin{tabular}{|c|p{4.8cm}|p{5.4cm}|}
\hline
测试对象 & 测试内容 & 预期结果 \\\hline
认证服务 & 登录认证、令牌生成、令牌解析 & 用户身份恢复正确，非法令牌被拒绝 \\\hline
订单服务 & 订单创建、审核、状态推进 & 状态按既定流程流转，不允许越级跳转 \\\hline
库存服务 & 入库、出库、库存流水更新 & 当前库存与流水记录保持一致 \\\hline
生产服务 & 生产计划生成、任务完成后转待检 & 生产任务状态更新正确，自动触发质检流程 \\\hline
采购服务 & 采购申请、采购单生成、到货处理 & 采购状态与供应商反馈保持同步 \\\hline
计划服务 & 库存预警判断、周计划计算 & 能识别供需缺口并生成合理建议 \\\hline
\end{tabular}
\end{table}

单元测试结果表明，系统关键服务层逻辑能够按照设计规则稳定运行，尤其是在状态控制和权限校验方面能够较好地保证业务一致性。

\section{功能测试}

功能测试以实际业务场景为中心，重点验证多角色协同流程是否通畅。测试覆盖登录认证、顾客下单、销售审核、库存查询、仓库发货、库存不足触发生产、生产完成进入质检、采购协同、供应商发货、质量追溯查询和实时消息通知等功能。功能测试结果如表~\ref{tab:function-test}所示。

\begin{table}[htbp]
\centering
\footnotesize
\caption{主要功能测试结果}
\label{tab:function-test}
\begin{tabular}{|c|p{5.2cm}|c|}
\hline
测试场景 & 测试要点 & 结果 \\\hline
登录认证 & 不同角色登录后进入对应功能范围，非法访问被拦截 & 通过 \\\hline
订单处理 & 顾客下单后进入销售审核，审核结果可回传 & 通过 \\\hline
库存管理 & 出入库操作后库存数量与流水同步更新 & 通过 \\\hline
生产流转 & 库存不足时触发生产，生产完成后进入待质检 & 通过 \\\hline
采购协同 & 采购订单可下发给供应商，供应商可反馈发货状态 & 通过 \\\hline
质量追溯 & 可依据订单号查询生产与质检链路信息 & 通过 \\\hline
实时通知 & 审核、预警、计划与质检结果可实时推送 & 通过 \\\hline
周计划生成 & 可生成生产与采购建议并通知相关角色 & 通过 \\\hline
\end{tabular}
\end{table}

从功能测试情况看，系统能够支撑主要业务场景的连续处理，不同角色之间的协作关系在系统中得到了较为清晰的体现。

\section{测试结果分析}

综合单元测试与功能测试结果可以看出，本文设计并实现的企业供应链协同管理系统能够基本满足预期设计要求。系统在认证授权、订单处理、库存流转、生产协同、采购协同、质检留痕、质量追溯和实时通知等方面均表现出较好的功能完整性。尤其是在多角色协同闭环和状态一致性方面，系统通过统一的服务层规则控制和消息同步机制，有效降低了人工流转过程中信息延迟或遗漏的风险。

但从测试结果来看，系统仍存在进一步优化空间。首先，周计划生成目前主要基于规则计算，对需求波动的适应能力仍有限；其次，实时通知机制已能满足一般业务协同需要，但在更大并发量下仍需进一步评估性能表现；再次，当前系统主要面向 Web 端访问，移动端适配能力尚有待增强。总体而言，系统已具备较好的工程可用性，但在计划智能化、性能优化和终端扩展方面仍有继续提升的空间。